% !TeX root = ../papers/nc-p3.tex
\section{交代形式と自己リンケージによる四次曲面の排除}
\label{sec:cofactor-integral-octic}
\begin{theorem}\label{thm:cofactor-integral-octic}
被約四辺有限窓の一般可逆な全階数反対根で、余因子スキームの曲線部分全体が一般被約で
その台 $C$ が整なら、四次以下の曲面は $C$ を含まない。像二平面のPlücker曲線は
$\mathbb P^5$ を張り、正規化上のPlücker次数 $\beta$ は五以上である。
八次行列式は既約であり、$17\le\deg C\le21$ である。
\end{theorem}
\begin{proof}[交代な縮約と仮定した四次曲面上の残余]
まず定理\ref{thm:cofactor-quartic-contact}までの結果だけを使う。
像二平面がある定数交代形式 $B\ne0$ について全て等方的なら、
$N(x)BN(x)$ の全成分は $C$ 上零な四次式である。それらが全て零なら一般可逆性により $B=0$。
そうでなければ $C$ を含む四次式は存在し、前定理によりその空間は一つの四次式 $F$ の定数倍である。
ところが $F$ の一般点では $N$ は可逆であり、同じ恒等式で再び $B=0$ となる。
これは像二平面のPlücker座標に線形関係がないことを示す。
正規化の線束には独立な六切断があるので次数 $\beta\ge5$ である。

$F$ が存在すると仮定すると、前定理で $d=16$ かつ $(F,\det N)$ は $C$ を台に持つ二倍曲線である。
この完全交叉 $Y$ 内での $C$ の残余は $C$ 自身である。
実際一般点では $F$ は滑らかで $\det N|_F$ は二重なので残余は一般点で単純、
完全交叉内のCohen--Macaulay曲線の残余は純曲線である。台と一般点での一致から両者は等しい。
したがって $\mathcal I_{C/Y}=\mathcal Hom_Y(\mathcal O_C,\mathcal O_Y)
=\omega_C(-8)$（$\omega_Y=\mathcal O_Y(8)$）。
四次で $H^1(\mathcal I_Y(4))=0$、かつ $I_C(4)=I_Y(4)=\langle F\rangle$ なので
$H^0(\omega_C(-4))=0$。Riemann--Rochと四次式の制限から
$65-p_a(C)=h^0(\mathcal O_C(4))\ge34$、従って $p_a(C)\le31$。
一方命題\ref{prop:cofactor-singular-length}は
$40\ge2m\ge96+2-2p_a(C)+\beta\ge41$ を与える。これで四次曲面を排除した。
\end{proof}

\begin{proof}[曲線次数と判別式の整性]
$p_a=p_a(C)$ と書く。同じ印数評価より $2p_a\ge6d+2+5-40$、
すなわち $p_a\ge3d-16$。四次式三十五個は $C$ 上で独立である。
$H^1(\mathcal O_C(4))=0$ なら $35\le4d+1-p_a\le d+17$ で $d\ge18$。
特殊なら整曲線のClifford不等式 $35\le1+2d$ で $d\ge17$ である。
ここで特異曲線でも、$L=\mathcal O_C(4)$ の切断と $\omega_C L^{-1}$ の切断の積は
どちらの因子も非零なら非零であり、積の包の次元は
$h^0(L)+h^1(L)-1$ 以上である。実際滑らかな一点で付値の異なる基底を両方に取り、
二つの整数集合の和集合（各一元の和の全体）の大きさを使えばよい。
その積は $H^0(\omega_C)$ に入り、Riemann--Rochを適用すると上の不等式が出る。

前定理の一意な二重軌跡を持つ既約因子を $G$、その次数を $e$ とする。
一般平面は $G$ を整な平面曲線で切り、相異なる $d$ 点で特異である。
各特異点の正規化余長は一以上だから $d\le(e-1)(e-2)/2$ である。
$e\le7$ なら右辺は十五以下となり不可能。従って $e=8$、行列式自体が既約で $d\le21$ となる。
\end{proof}
